\documentclass[10pt]{amsart}
\usepackage{amsmath,amssymb,amsthm}
\setlength{\textwidth}{\paperwidth}\addtolength{\textwidth}{-3.6cm}
\setlength{\oddsidemargin}{-0.74cm}\setlength{\evensidemargin}{-0.74cm}
\setlength{\headheight}{0pt}\setlength{\headsep}{0pt}
\setlength{\topmargin}{-0.84cm}
\setlength{\textheight}{\paperheight}\addtolength{\textheight}{-3.6cm}
\setlength{\footskip}{0.8cm}
\ifdefined\pdfpagewidth\pdfpagewidth=\paperwidth\pdfpageheight=\paperheight\fi
\pagestyle{plain}\raggedbottom
\setlength{\parskip}{0pt plus 0.5pt}
\widowpenalty=150 \clubpenalty=150
\makeatletter\def\section{\@startsection{section}{1}{\z@}{.6\baselineskip}{.25\baselineskip}{\normalfont\bfseries}}\makeatother
\newtheorem{thm}{Theorem}
\newtheorem{lem}[thm]{Lemma}
\newcommand{\st}[1]{\par\noindent\textit{#1}\enspace}

\title{The $\pm1$ case from a one-parameter family}
\date{}
\begin{document}\sloppy
\maketitle
\vspace{-2.5em}

\noindent The argument is DLME's (arXiv:2410.21248, Props.~3.9--3.10 and \S5). It is
restated here so that the role of each ingredient is visible, and read in the
pillowcase language of the uploaded paper.

\section{Statement}

For a homology sphere $Y$, $I_\bullet(Y)$ is the $\mathbb Z$-graded, Chern--Simons
filtered irreducible complex. $\kappa(d)$ is the least filtration level of a
nonzero class in degree $d$, and $\ell(Y)=\inf_d(\kappa(d)-d/8)$. In the
normalisation used throughout, this is the least $(3+\rho)/16$ of a nonzero class.

\begin{thm}[DLME]\label{t}
Let $K\subset S^3$ be nontrivial. Then $\ell(S^3_{-1}(K))\le\ell(S^3_1(K))-\frac18-\eta(K)$ with
$\eta(K)>0$. In particular $S^3_1(K)\not\cong S^3_{-1}(K)$ as oriented manifolds.
\end{thm}

\section{Ingredients}

\st{(1) The triangle.} For $K\subset Y$ there is a distance-two exact triangle
(CDX; DLME Prop.~3.9, over $\mathbb F_2$)
\[
\cdots\to I(Y_1)\xrightarrow{\;W_*\oplus(W,c)_*\;}I(Y)\oplus I(Y)\xrightarrow{\;(W',c')_*\oplus W'_*\;}I(Y_{-1})\to\cdots .
\]
\emph{Pillowcase reading} (heuristic: it uses the paper's Conjecture~3.2). $L_{-1}$ is obtained from $L_1$ by one Dehn twist along
each edge $\{a=0\}$ and $\{a=\frac12\}$. Both edges are the line of $Y=S^3$: they are
$L_{1/0}$ and its $\chi$-translate, which is why $I(Y)$ appears twice. Each twist
contributes one Seidel triangle whose third term is the edge. $\Gamma_K$ misses
both edges (Prop.~1.1(b)), so for $Y=S^3$ the third term contributes nothing.
This is the distance-two analogue of the uploaded paper's \S4.2.

\st{(2) The one-parameter family map.} Triangle detection produces a null-homotopy
$g_1\colon I_\bullet(Y_1)\to I_\bullet(Y_{-1})$ of the composite through the middle
term. It counts instantons with irreducible ends over a one-parameter family of
metrics $G\cong[-1,1]$ on the composite cobordism $W^1_{-1}=W(Y_1\to S^3)\cup_{S^3}W(S^3\to Y_{-1})$.
At one end the family breaks along the middle $S^3$. At the other it breaks along a
hypersurface with $\mathbb RP^3$ components (DLME \S5). The $\mathbb RP^3$ bounds the $(-2)$-sphere
formed by the cocore of the first handle and the core of the second. When the middle term is $0$, as for $Y=S^3$, $g_1$ is a
chain map and a quasi-isomorphism (DLME Prop.~3.10, \S5).

\st{(3) Its level.} $W^1_{-1}$ has $b_1=b^+=0$ and $\pi_1=1$: each piece is a
$2$-handle attached to $S^3\times I$. An instanton counted over a $d$-parameter
family has index $i_0=-d$. The level formula
$c(\beta)-c(\alpha)=\frac{i_0+3b^+}8-E$ then gives
\[
c(\beta)-c(\alpha)=-\tfrac18-E .
\]
The bundle drops out. In DLME's terms, with $c^2=-1$: degree $D=1-2c^2=3$, level
$L=-c^2/4-E$, and $L-D/8=-\frac18-E$.

\st{(4) Positivity.} $E>0$, because a flat connection on the simply connected $W^1_{-1}$
is trivial and cannot have irreducible limits. Compactness gives finitely many
counted instantons, so $\eta(K):=\min E>0$.

\st{(5) Nonvanishing.} $I(S^3_{\pm1}(K))\cong I^w(S^3_0(K))\ne0$ by Floer's triangle and
Kronheimer--Mrowka. Hence $\ell$ is finite.

\begin{lem}[DLME Lemma 2.7]
Let $f$ be a filtered chain map of degree $D$ and level $L$ that is injective on
homology in every degree. Then $\ell(\text{target})\le\ell(\text{source})+L-D/8$.
\end{lem}
\noindent\emph{Proof.} A class at level $\kappa(d)+\epsilon$ maps to a nonzero class at
level $\kappa(d)+L+\epsilon$ in degree $d+D$. Subtract $(d+D)/8$ and take infima. $\square$

\noindent\emph{Proof of Theorem \ref{t}.} Apply the lemma to $g_1$, using (2), (3), (4)
and (5). $\square$

\section{What each ingredient contributes}

\begin{itemize}
\item \emph{The $\frac18$ is the family dimension, and nothing else.} One parameter means
one unit of index, which is $\frac18$ of level. No $\rho$-term from a middle end
enters the counted maps. The $\mathbb RP^3$ appears only at the boundary of the
family.
\item \emph{Strictness is positivity of energy} on a simply connected cobordism, which
is an allowed input.
\item \emph{What the cones of the uploaded paper's \S6 cannot supply} is the direction
of the inequality. Those cones join $Y_{\pm1}$ through $(Y_0,w)$; that is the
horizontal rotation, and it gives only
$\ell(Y_1)-\ell(Y_{-1})=\frac38-\operatorname{Area}(T)-\frac18\Delta\iota$. For
a pair $x\mapsto y$ counted by $g_1$, Stokes and the level formula (exact for torus
knots, plausible in general) give
\[
E=\tfrac14-\operatorname{Area}(T_{xy})-\tfrac18\Delta\iota .
\]
So $E>0$ is precisely the area bound that the cones cannot see. For the
right-handed trefoil, $\ell(Y_1)-\ell(Y_{-1})=\frac{61}{280}=\frac18+\frac{13}{140}$.
\end{itemize}

\section{Why the same proof does not give $\pm2$}

$L_{-2}$ is obtained from $L_2$ by \emph{two} twists along each edge. The relevant
composite $W'=(-X_2)\cup_{S^3}X_{-2}$ contains a $(-4)$-sphere bounded by $L(4,1)$. The
analogous family map has the same index count, so the same level $-\frac18-E$, and $E>0$
holds again because $W'$ is simply connected. What fails is one of the two
structural properties:
\begin{itemize}
\item \emph{Bundle even on the sphere:} flat caps on the $(-4)$-disk bundle with central
$L(4,1)$ boundary give uncancelled lens faces. The middle term is not $0$, and the
family map is not a chain map. This matches the uploaded paper's \S4.5 remark that $(\frac12,0)$
is a central flat connection on $Y_{\pm2}$; the identification is my reading.
\item \emph{Bundle odd on the sphere (mixed bundles):} the family map is plausibly a chain
map, but it is not injective in general. It would force $m\ge\frac18$, whereas
$m_2=\frac1{16}$ for $T(2,-3)$: the class is absorbed.
\end{itemize}
These are DLME's two ``inadequate variations''. The map that is a quasi-isomorphism,
$B$, has level $+\frac18$, because the insertion $\mu(S)$ costs $\frac14$. So the
$\pm2$ problem is exactly to find a replacement for $g_1$ that is injective with
level below $0$. The $\pm1$ argument needs nothing beyond the family dimension and
positivity because there the middle term vanishes on the nose.
\end{document}
